\clearpage
\section{Especificação formal M13: Notificação unificada}
\label{sec:formal-m13}
\sloppy

M13 formaliza a entidade única \textit{Notificacao} e a caixa (UI-12) como uma
\textbf{composição} das fatias já validadas: idempotência M7, escassez e docId
diário M8, autorização M9, participação canônica M11, efeitos delimitados de
M12.1/M12.2. O contrato estrutural está em
\texttt{specification/cue/domain/formal\_m13.cue}
(\texttt{\#M13Contrato}) e a prova relacional/temporal em
\texttt{specification/alloy/operations/notificacoes\_m13.als}. O modelo é
composto, não justaposto: caixa, leitura, marcação, ``Limpar tudo'', expiração,
alvo e emissão compartilham o mesmo estado, e os predicados de autorização
(\texttt{versaoCorrente}, \texttt{authOk}, \texttt{temVinculo}) são reproduzidos
de M12.1/M9 com guard de drift (\texttt{tools/formal/m13\_composition.mjs}). A
fonte normativa humana é a Seção~\ref{sec:regras-notificacoes-m13}.

\subsection*{Caixa única, papel visual e Bolsista/Q12}
A visibilidade normativa é a posse (\texttt{caixaDe[s,u,n]}), independente do
papel selecionado no header; o papel armazenado é contexto de apresentação e
nunca concede recurso. Um usuário não lê a caixa de outro UID
(\texttt{NaoLeCaixaAlheia}); o papel visual não filtra a própria caixa nem
concede acesso ao alvo (\texttt{PapelVisualNaoFiltraCaixa},
\texttt{PapelVisualNaoConcedeRecurso}); o Bolsista vê os alertas próprios
(\texttt{BolsistaVeCaixa}, \texttt{WitnessBolsistaVeAlerta}) e uma conta
multi-role reúne avisos de papéis distintos
(\texttt{WitnessCaixaMultiRole}).

\subsection*{Marcação e ``Limpar tudo''}
Marcar lida é restrito ao dono, idempotente e não reverte
(\texttt{MarcarSoProprias}, \texttt{MarcaIdempotente}, \texttt{MarcaNaoReverte});
``Limpar tudo'' marca apenas itens do próprio UID, por corte, sem
\texttt{DELETE}, e um aviso emitido após o corte permanece para a próxima rodada
(\texttt{LoteSoProprias}, \texttt{SemDelete}, \texttt{LimparTudoCorteEstavel},
\texttt{WitnessLimparTudoCorteComNovaEmissao}); o retry retoma do corte sem
deixar itens presos (\texttt{WitnessLimparTudoRetry}).

\subsection*{Expiração e alvo}
\texttt{expira\_em = null} não expira automaticamente
(\texttt{ExpiracaoDistingueNull}) e a expiração retira do conjunto ativo sem
apagar o fato (\texttt{ExpiradoForaDoAtivo}); \texttt{ESCASSEZ\_ESTOQUE}
permanece nulo e persistente (\texttt{WitnessNullNaoExpira}). \texttt{entidade\_alvo}
é rota conhecida e \texttt{id\_alvo} é identificador de recurso; URL arbitrária
não navega (\texttt{AlvoInvalidoNaoNavega}, \texttt{UrlArbitrariaNaoNavega}).
Abrir revalida autorização corrente: sem vínculo canônico, com objeto removido ou
sem escopo Q13 do Chefe, não há acesso
(\texttt{VinculoCanonicoRequerido}, \texttt{ObjetoRemovidoNaoNavega},
\texttt{ChefeSemEscopoNaoAcessa}); a revogação corta o clique sem apagar o fato
(\texttt{RevogacaoCortaClique}, \texttt{WitnessRevogacaoPreservaFato}); a
notificação nunca contorna a ACL (\texttt{AlertaNaoContornaAcl}). As pontes
\texttt{PonteM9ExigeAuth} e \texttt{PonteM12NaoAfrouxaAcl} provam que o acesso
derivado exige autorização e não é ampliado pela caixa.

\subsection*{Privacidade e emissão/deduplicação}
O payload é mínimo: nenhuma notificação carrega conteúdo protegido
(\texttt{NotificacaoNaoExpoeOriginal}) e Firestore Rules não mascaram campos.
A emissão reutiliza a identidade M7 e a chave determinística M8: a mesma
identidade não duplica (\texttt{EmissaoMesmaIdentidadeNaoDuplica},
\texttt{ComposicaoM7RetryNaoDuplica}), chaves distintas podem emitir
(\texttt{EmissaoChaveDistintaEmite}) e erro de emissão não vira sucesso
(\texttt{ErroEmissaoNaoViraSucesso}). O contexto é condicional: \texttt{id\_turma}
obrigatório em tipo acadêmico e nulo nos demais
(\texttt{IdTurmaAcademicoObrigatorio}, \texttt{IdTurmaOperacionalNulo}).

\subsection*{Rastreabilidade norma $\rightarrow$ CUE/IR $\rightarrow$ prova $\rightarrow$ receipt}
\begingroup\scriptsize
\setlength{\tabcolsep}{3pt}
\begin{longtable}{>{\raggedright\arraybackslash}p{0.20\textwidth} >{\raggedright\arraybackslash}p{0.18\textwidth} >{\raggedright\arraybackslash}p{0.40\textwidth} >{\raggedright\arraybackslash}p{0.16\textwidth}}
\toprule
\textbf{Norma (Seção 7.8)} & \textbf{CUE/IR} & \textbf{Check/Witness (Alloy)} & \textbf{Receipt} \\
\midrule
\endfirsthead
\toprule
\textbf{Norma} & \textbf{CUE/IR} & \textbf{Prova} & \textbf{Receipt} \\
\midrule
\endhead
\bottomrule
\endfoot
RN-M13-01 & \texttt{\#M13Notificacao}, \texttt{\#M13Caixa} & \texttt{NaoLeCaixaAlheia}, \newline \texttt{PapelVisualNaoFiltraCaixa}, \newline \texttt{BolsistaVeCaixa}, \newline \texttt{WitnessCaixaMultiRole} & M13-INV-002..005, M13-WIT-029/030 \\
RN-M13-02 & \texttt{\#M13Marcacao}, \texttt{\#M13LimparTudo} & \texttt{MarcarSoProprias}, \newline \texttt{MarcaIdempotente}, \newline \texttt{LoteSoProprias}, \newline \texttt{SemDelete}, \newline \texttt{LimparTudoCorteEstavel} & M13-INV-006..011, M13-WIT-031..035 \\
RN-M13-03 & \texttt{\#M13Expiracao} & \texttt{ExpiracaoDistingueNull}, \newline \texttt{ExpiradoForaDoAtivo}, \newline \texttt{WitnessNullNaoExpira} & M13-INV-012/013, M13-WIT-036/037 \\
RN-M13-04 & \texttt{\#M13Alvo} & \texttt{AlvoInvalidoNaoNavega}, \newline \texttt{UrlArbitrariaNaoNavega}, \newline \texttt{AlertaNaoContornaAcl}, \newline \texttt{RevogacaoCortaClique}, \newline \texttt{ChefeSemEscopoNaoAcessa}, \newline \texttt{PonteM9ExigeAuth} & M13-INV-014..020, \newline M13-INV-027/028, \newline M13-WIT-038..040/044 \\
RN-M13-05 & \texttt{contem\_conteudo\_protegido} & \texttt{NotificacaoNaoExpoeOriginal} & M13-INV-021 \\
RN-M13-06 & \texttt{\#M13Emissao}, \texttt{\#M13Operacao} & \texttt{EmissaoMesmaIdentidadeNaoDuplica}, \newline \texttt{ComposicaoM7RetryNaoDuplica}, \newline \texttt{EmissaoChaveDistintaEmite}, \newline \texttt{ErroEmissaoNaoViraSucesso} & M13-INV-022..024, \newline M13-WIT-041..043 \\
RN-M13-07 & \texttt{id\_turma} condicional & \texttt{IdTurmaAcademicoObrigatorio}, \newline \texttt{IdTurmaOperacionalNulo} & M13-INV-025/026, \newline M13-WIT-045/046 \\
RN-M13-08 & reuso M12.2 (\texttt{geracao}) & \texttt{AnexoCompostoUsaGeracaoCanonica} (M12) & receipt M12 \\
\bottomrule
\end{longtable}
\endgroup

\begingroup\small\raggedright
% Gerado por lcqui-spec-doc; não editar.
\subsection*{Evidência formal M13 --- Notificação unificada}
CUE verifica os shapes de Notificação unificada (caixa, leitura, marcação, Limpar tudo, alvo, expiração, emissão idempotente e operação), com enums fechados de tipo/papel/alvo, condicionais (lida/lida\_em, id\_turma obrigatório nos seis tipos acadêmicos e nulo nos demais, payload mínimo, expiração nula em ESCASSEZ\_ESTOQUE) e limites. Alloy compõe a caixa única por UID (um único sino que reúne todos os papéis, sem o papel visual filtrar nem conceder privilégio), a leitura restrita ao destinatário, a marcação individual e o Limpar tudo paginado/reentrante com corte estável e sem DELETE, a expiração (NULL não expira; expirado sai do ativo sem apagar) e a privacidade do payload. A autorização de Posts, Comentários, Roteiros, compartilhamento e escopo Q13 é reproduzida de composition\_m12.als (M12.1/M12.2/M9/M11) com guard de drift: a rota acadêmica exige papel acadêmico e vínculo canônico atual (Chefe com vínculo legado não a usa), o professor dono e o compartilhamento preservam seus acessos, e a emissão de URL de Roteiro por Chefe exige escopo de auditoria M12.2. A emissão/deduplicação reutiliza a identidade M7 e a chave determinística M8. Rust valida proveniência, ordem exata do receipt e as origens de composição. A evidência não certifica Firebase, Firestore/Storage, Rules, Auth, relógio real, paginação/concorrência real nem a aplicação.
\begin{description}
\item[M13-INV-001] TransicoesPreservamCoerencia: UNSAT.\newline Escopo: \texttt{check TransicoesPreservamCoerencia for 4}.
\item[M13-INV-002] NaoLeCaixaAlheia: UNSAT.\newline Escopo: \texttt{check NaoLeCaixaAlheia for 4}.
\item[M13-INV-003] PapelVisualNaoFiltraCaixa: UNSAT.\newline Escopo: \texttt{check PapelVisualNaoFiltraCaixa for 4}.
\item[M13-INV-004] BolsistaVeCaixa: UNSAT.\newline Escopo: \texttt{check BolsistaVeCaixa for 4}.
\item[M13-INV-005] MarcarSoProprias: UNSAT.\newline Escopo: \texttt{check MarcarSoProprias for 4}.
\item[M13-INV-006] MarcaIdempotente: UNSAT.\newline Escopo: \texttt{check MarcaIdempotente for 4}.
\item[M13-INV-007] MarcaNaoReverte: UNSAT.\newline Escopo: \texttt{check MarcaNaoReverte for 4}.
\item[M13-INV-008] LoteSoProprias: UNSAT.\newline Escopo: \texttt{check LoteSoProprias for 4}.
\item[M13-INV-009] SemDelete: UNSAT.\newline Escopo: \texttt{check SemDelete for 4}.
\item[M13-INV-010] LimparTudoCorteEstavel: UNSAT.\newline Escopo: \texttt{check LimparTudoCorteEstavel for 5}.
\item[M13-INV-011] ExpiracaoDistingueNull: UNSAT.\newline Escopo: \texttt{check ExpiracaoDistingueNull for 4}.
\item[M13-INV-012] ExpiradoForaDoAtivo: UNSAT.\newline Escopo: \texttt{check ExpiradoForaDoAtivo for 4}.
\item[M13-INV-013] AlvoInvalidoNaoNavega: UNSAT.\newline Escopo: \texttt{check AlvoInvalidoNaoNavega for 4}.
\item[M13-INV-014] AlertaNaoContornaAcl: UNSAT.\newline Escopo: \texttt{check AlertaNaoContornaAcl for 4}.
\item[M13-INV-015] RotaAcademicaExigePapel: UNSAT.\newline Escopo: \texttt{check RotaAcademicaExigePapel for 4}.
\item[M13-INV-016] ChefeLegadoNaoUsaAcademico: UNSAT.\newline Escopo: \texttt{check ChefeLegadoNaoUsaAcademico for 4}.
\item[M13-INV-017] SemVinculoNaoRestaura: UNSAT.\newline Escopo: \texttt{check SemVinculoNaoRestaura for 4}.
\item[M13-INV-018] ProfessorDonoMantemAcesso: UNSAT.\newline Escopo: \texttt{check ProfessorDonoMantemAcesso for 4}.
\item[M13-INV-019] CompartilhamentoSoProfessor: UNSAT.\newline Escopo: \texttt{check CompartilhamentoSoProfessor for 4}.
\item[M13-INV-020] ChefeAdminNaoHerdaAcademico: UNSAT.\newline Escopo: \texttt{check ChefeAdminNaoHerdaAcademico for 4}.
\item[M13-INV-021] ChefeSemEscopoNaoEmiteUrl: UNSAT.\newline Escopo: \texttt{check ChefeSemEscopoNaoEmiteUrl for 4}.
\item[M13-INV-022] EscopoEncerradoImpedeNovaEmissao: UNSAT.\newline Escopo: \texttt{check EscopoEncerradoImpedeNovaEmissao for 4}.
\item[M13-INV-023] EncerramentoEscopoPreservaRoteiro: UNSAT.\newline Escopo: \texttt{check EncerramentoEscopoPreservaRoteiro for 4}.
\item[M13-INV-024] RoteiroAlunoExigePostEAnexo: UNSAT.\newline Escopo: \texttt{check RoteiroAlunoExigePostEAnexo for 4}.
\item[M13-INV-025] PonteRotaChefeRoteiroRefinaUrl: UNSAT.\newline Escopo: \texttt{check PonteRotaChefeRoteiroRefinaUrl for 4}.
\item[M13-INV-026] NotificacaoNaoExpoeOriginal: UNSAT.\newline Escopo: \texttt{check NotificacaoNaoExpoeOriginal for 4}.
\item[M13-INV-027] ColegaVeAviso: UNSAT.\newline Escopo: \texttt{check ColegaVeAviso for 4}.
\item[M13-INV-028] AutorVeOriginal: UNSAT.\newline Escopo: \texttt{check AutorVeOriginal for 4}.
\item[M13-INV-029] AuditorVeOriginal: UNSAT.\newline Escopo: \texttt{check AuditorVeOriginal for 4}.
\item[M13-INV-030] EmissaoMesmaIdentidadeNaoDuplica: UNSAT.\newline Escopo: \texttt{check EmissaoMesmaIdentidadeNaoDuplica for 4}.
\item[M13-INV-031] ErroEmissaoNaoViraSucesso: UNSAT.\newline Escopo: \texttt{check ErroEmissaoNaoViraSucesso for 4}.
\item[M13-INV-032] ComposicaoM7RetryNaoDuplica: UNSAT.\newline Escopo: \texttt{check ComposicaoM7RetryNaoDuplica for 4}.
\item[M13-INV-033] IdTurmaAcademicoObrigatorio: UNSAT.\newline Escopo: \texttt{check IdTurmaAcademicoObrigatorio for 4}.
\item[M13-INV-034] IdTurmaOperacionalNulo: UNSAT.\newline Escopo: \texttt{check IdTurmaOperacionalNulo for 4}.
\item[M13-WIT-035] WitnessCaixaMultiRole: SAT.\newline Escopo: \texttt{run WitnessCaixaMultiRole for 4}.
\item[M13-WIT-036] WitnessBolsistaVeAlerta: SAT.\newline Escopo: \texttt{run WitnessBolsistaVeAlerta for 4}.
\item[M13-WIT-037] WitnessMarcarLidaPropria: SAT.\newline Escopo: \texttt{run WitnessMarcarLidaPropria for 4}.
\item[M13-WIT-038] WitnessMarcacaoIdempotente: SAT.\newline Escopo: \texttt{run WitnessMarcacaoIdempotente for 4}.
\item[M13-WIT-039] WitnessLimparTudoProprias: SAT.\newline Escopo: \texttt{run WitnessLimparTudoProprias for 4}.
\item[M13-WIT-040] WitnessLimparTudoRetry: SAT.\newline Escopo: \texttt{run WitnessLimparTudoRetry for 5}.
\item[M13-WIT-041] WitnessLimparTudoCorteComNovaEmissao: SAT.\newline Escopo: \texttt{run WitnessLimparTudoCorteComNovaEmissao for 5}.
\item[M13-WIT-042] WitnessExpiraAtivo: SAT.\newline Escopo: \texttt{run WitnessExpiraAtivo for 4}.
\item[M13-WIT-043] WitnessNullNaoExpira: SAT.\newline Escopo: \texttt{run WitnessNullNaoExpira for 4}.
\item[M13-WIT-044] WitnessAlvoInvalidoNaoNavega: SAT.\newline Escopo: \texttt{run WitnessAlvoInvalidoNaoNavega for 4}.
\item[M13-WIT-045] WitnessRotaAcademicaValida: SAT.\newline Escopo: \texttt{run WitnessRotaAcademicaValida for 4}.
\item[M13-WIT-046] WitnessArquivamentoTurma: SAT.\newline Escopo: \texttt{run WitnessArquivamentoTurma for 4}.
\item[M13-WIT-047] WitnessDesarquivamentoTurma: SAT.\newline Escopo: \texttt{run WitnessDesarquivamentoTurma for 4}.
\item[M13-WIT-048] WitnessChefeLegadoNaoUsaAcademico: SAT.\newline Escopo: \texttt{run WitnessChefeLegadoNaoUsaAcademico for 4}.
\item[M13-WIT-049] WitnessSemVinculoClaimAtual: SAT.\newline Escopo: \texttt{run WitnessSemVinculoClaimAtual for 4}.
\item[M13-WIT-050] WitnessProfessorDonoPost: SAT.\newline Escopo: \texttt{run WitnessProfessorDonoPost for 4}.
\item[M13-WIT-051] WitnessCompartilhamentoRoteiro: SAT.\newline Escopo: \texttt{run WitnessCompartilhamentoRoteiro for 4}.
\item[M13-WIT-052] WitnessChefeComEscopo: SAT.\newline Escopo: \texttt{run WitnessChefeComEscopo for 4}.
\item[M13-WIT-053] WitnessEscopoEncerrado: SAT.\newline Escopo: \texttt{run WitnessEscopoEncerrado for 4}.
\item[M13-WIT-054] WitnessEmissaoUnica: SAT.\newline Escopo: \texttt{run WitnessEmissaoUnica for 4}.
\item[M13-WIT-055] EmissaoChaveDistintaEmite: SAT.\newline Escopo: \texttt{run EmissaoChaveDistintaEmite for 4}.
\item[M13-WIT-056] WitnessErroNaoEmite: SAT.\newline Escopo: \texttt{run WitnessErroNaoEmite for 4}.
\item[M13-WIT-057] WitnessColegaVeAviso: SAT.\newline Escopo: \texttt{run WitnessColegaVeAviso for 4}.
\item[M13-WIT-058] WitnessAutorVeOriginal: SAT.\newline Escopo: \texttt{run WitnessAutorVeOriginal for 4}.
\item[M13-WIT-059] WitnessAuditorVeOriginal: SAT.\newline Escopo: \texttt{run WitnessAuditorVeOriginal for 4}.
\end{description}
UNSAT para check indica ausência de contraexemplo no escopo declarado; SAT para run indica testemunha encontrada. A evidência não certifica a implementação Firebase/backend.

\par\endgroup

\subsection*{Limites}
Escopos limitados (\texttt{for 4}/\texttt{for 5}). A prova composta não
certifica \texttt{functions/}, \texttt{frontend/}, Firestore/Storage Rules, Auth,
Storage real, entrega de e-mail, relógio de produção (fuso
América/São Paulo), paginação e concorrência real do Firestore nem atomicidade
entre serviços. O receipt registra \texttt{model\_sha256},
\texttt{spec\_ir\_sha256} e as origens de composição (M7/M8/M9/M12). O capítulo
não reabre nem substitui as provas isoladas dessas fatias.
\fussy
